\documentclass[../../../include/open-logic-section]{subfiles}

\begin{document}

\olfileid{sth}{ord-arithmetic}{expo}
\olsection{توان‌رسانیِ ترتیبی}

اکنون به توان‌رسانیِ ترتیبی می‌پردازیم. متأسفانه برای آن تعریف
ترکیبیِ \emph{خوشایندی} وجود ندارد.

البته تعریف‌های ترکیبیِ صریح \emph{وجود دارند}. این یکی را در نظر
بگیرید. بگذارید $\text{finfun}(\alpha,\beta)$ مجموعهٔ همهٔ توابع
$f \colon \alpha \to \beta$ باشد به‌طوری‌که
$\Setabs{\gamma \in \alpha}{f(\gamma) \neq 0}$ با یک عدد طبیعی
هم‌شمار باشد. روی $\text{finfun}(\alpha,\beta)$ چنین خوش‌ترتیبی
تعریف کنید: $f \sqsubset g$ اگر و تنها اگر $f \neq g$ و
$f(\gamma_0) < g(\gamma_0)$، که در آن
$\gamma_0 = \text{max}\Setabs{\gamma \in \alpha}{f(\gamma) \neq
g(\gamma)}$. آن‌گاه می‌توانیم $\ordexpo{\alpha}{\beta}$ را به‌صورت
$\ordtype{\text{finfun}(\beta, \alpha), \sqsubset}$ تعریف کنیم.
پاتر این تعریف صریح را به کار می‌گیرد و بی‌درنگ توضیح می‌دهد:
\begin{quote}
	انتخاب این ترتیب صرفاً از خواست ما برای به‌دست‌آوردن تعریفی از
	توان‌رسانیِ ترتیبی تعیین می‌شود که شرط بازگشتیِ مناسب را برآورده
	کند\ldots، و تجسم آن بسیار دشوارتر از جمعِ مرتب یا ضربِ مرتب است.
	\citep[p.~199]{Potter2004}
\end{quote}
دقیقاً. جمع را چنین توضیح دادیم: «یک نسخه از $\alpha$ و سپس یک نسخه
از $\beta$»؛ و ضرب را چنین: «دنباله‌ای $\beta$تایی از نسخه‌های
$\alpha$». اما دربارهٔ $\text{finfun}(\gamma,\alpha)$ هیچ بیان کوتاه
و گویایی نداریم. بنابراین، تعریف توان‌رسانیِ ترتیبی را
\emph{صرفاً با} بازگشت ترامتناهی عرضه می‌کنیم؛ یعنی:

\begin{defn}\ollabel{ordexporecursion}
\begin{align*}
	\ordexpo{\alpha}{0} &= 1\\
	\ordexpo{\alpha}{\beta\ordplus 1} &=\ordexpo{\alpha}{\beta} \ordtimes \alpha\\
	\ordexpo{\alpha}{\beta} &= \bigcup_{\delta < \beta}\ordexpo{\alpha}{\delta}& & \text{هرگاه $\beta$ عدد ترتیبیِ حدی باشد}
\end{align*}
\end{defn}

اگر \emph{به‌عنوان} نظریه‌پرداز مجموعه‌ها کار می‌کردیم، شاید
می‌خواستیم برخی ویژگی‌های توان‌رسانی ترتیبی را بررسی کنیم. اما مطلب
چندانی برای افزودن نداریم، جز این واقعیتِ غافلگیرنکننده که توان‌رسانی
ترتیبی جابجایی‌پذیر نیست. پس
$\ordexpo{2}{\omega} =
\bigcup_{\delta < \omega}\ordexpo{2}{\delta} = \omega$، درحالی‌که
$\ordexpo{\omega}{2} = \omega \ordtimes \omega$. البته نباید
\emph{انتظار} داشته باشیم توان‌رسانی جابجایی‌پذیر باشد، زیرا روی
اعداد طبیعی نیز چنین نیست:
$\ordexpo{2}{3} = 8 < 9 = \ordexpo{3}{2}$.

\begin{prob}
با استفاده از استقرای ترامتناهی اثبات کنید که اگر
$\ordexpo{\alpha}{\beta} =
\ordtype{\text{finfun}(\beta, \alpha), \sqsubset}$ را تعریف کنیم،
معادلات بازگشتیِ
\olref[sth][ord-arithmetic][expo]{ordexporecursion} را به دست
می‌آوریم.
\end{prob}

\end{document}
